\section{Backup and Disaster Recovery}\label{sec:backup}

Backup execution is healthy and current; backup \emph{architecture} is the weak point.
Live evidence (E-08) shows 14 snapshots dated 2026-07-23 from the previous night's runs,
with the PBS datastore at 97\,GB used of 23\,TB. Every cluster guest is protected daily.
The problem is that all of those copies live on one node.

\subsection{Backup Posture}

\begin{tabularx}{\linewidth}{@{}p{3.0cm} p{2.6cm} p{2.4cm} X@{}}
\toprule
\rh{Asset} & \rh{Target} & \rh{Schedule / retention} & \rh{Notes} \\
\midrule
\endhead
LXC 101/102/106/107 & randy-pbs \code{.10.187} & 02:00 daily, 7d/4w & NPM, Vaultwarden, Homepage, Open WebUI \\
LXC 103 Grafana & randy-pbs & 02:00 daily & observability stack \\
LXC 105 Headscale, 108 pihole2 & randy-pbs & 02:00 daily & pihole2 mirrors CT103 config \\
VM 100 OPNsense & randy-pbs + config repo & 03:00 daily + 03:17 & config also age-encrypted to GitHub \\
VM 104 Wazuh & randy-pbs-10g \code{.30.187} & 03:00 daily & 10G path \\
VM 110 HomeAssistant & randy-pbs & 03:30 daily & HAOS \\
VM 201/202/203 RKE2 & randy-pbs & 04:00 daily & etcd not separately backed up \\
QuarkyLab workspace & PBS host backup & 05:30 daily & fernanda refreservation floor \\
pve1 CT103/CT104 & \sevhigh NONE & never & primary Pi-hole + old homepage unbacked (QL-R06) \\
\code{bulk} pool data & \sevhigh NONE & never & media + fernanda research, no second copy (QL-R02) \\
Ares \code{/home} & restic to \code{randy:/mnt/bulk/...} & daily 04:03 & encrypted, but lands on the DEGRADED bulk pool \\
Ares \code{\textasciitilde/.claude} & \code{randy:/\allowbreak datastore/\allowbreak ares-claude} & hourly & memory/transcripts, secrets excluded \\
\bottomrule
\end{tabularx}

\subsection{3-2-1 Analysis}

The estate is effectively 1-1-1: one PBS instance, on one ZFS pool, in one rack, on one
UPS bus (QL-R03, \sevhigh). Three specific exposures compound it:

\begin{itemize}
\item \textbf{pve1 guests are unprotected} (QL-R06 / R-03): the primary Pi-hole CT103 and
the retired CT104 have no vzdump job. This is partly mitigated because the secondary
Pi-hole CT108 is backed up and nebula-syncs the primary's config, so DNS state is
rebuildable.
\item \textbf{The \code{bulk} pool has no second copy} (QL-R02): it is neither a PBS
source nor a target, so the DUNE research data and media it holds exist only on the
currently degraded DS4246 shelf.
\item \textbf{The Ares restic repository lands on \code{bulk}} (\code{randy:/mnt/bulk/
backups/ares}), the same degraded pool, so the workstation backup shares the shelf's
fault domain.
\end{itemize}

Encryption is partial: the OPNsense config backup is age-encrypted and the Ares restic
repository is encrypted, but the PBS datastore itself is not encrypted at rest.

\begin{callout}[title=Do not assume a green job is a recoverable backup]
PBS reports success and runs a weekly verify plus garbage collection, which is good, but
no full restore has been rehearsed for most guests. A backup that has never been restored
is an untested backup.
\end{callout}

\subsection{Recovery Priority and Objectives}

Recovery priority, highest first: (1) power and core network, (2) OPNsense routing and
DHCP, (3) DNS, (4) Randy storage and PBS, (5) Vaultwarden and certificates, (6) the app
and AI tiers, (7) research and student services. Estimated objectives from the backup
matrix: RPO is 24\,h for guest images (daily vzdump) and 1\,h for Ares memory; RTO ranges
from roughly 10-30 minutes per guest restore to several hours for a full cluster rebuild.
The OPNsense config RPO is effectively 24\,h with a 15-30 minute restore.

\subsection{Non-Destructive Restore Validation Plan}

No restore was performed during this review. The recommended validation, all
non-destructive and to be scheduled with operator approval, is: (a) run a PBS
\code{verify} job and confirm zero failures; (b) restore one representative guest to a
scratch VMID on a node with free space, boot it isolated, confirm it starts, then destroy
the scratch guest; (c) test-decrypt the OPNsense age config and the break-glass snapshot.
Never restore into a production VMID.

\subsection{Per-Domain Recovery Procedures}

The estate has an unusually strong recovery-documentation culture; the procedures below
already exist and were reviewed:

\begin{itemize}
\item \textbf{OPNsense:} serial console via \code{qm terminal 100} from pve2 plus a
DR-tested age-encrypted \code{config.xml} restore; SSH is disabled by design so console
is the only shell.
\item \textbf{DNS:} \code{pct start 103} on pve1, or fail to the secondary \code{.178};
\code{pihole restartdns} if the container is up but DNS is stuck.
\item \textbf{Cluster/secrets loss:} the break-glass age snapshot held on both Ares and
Randy breaks the Vaultwarden-in-cluster circular dependency (QL-R08), letting an operator
read credentials with Vaultwarden, NPM, and the cluster all down.
\item \textbf{Network device:} the EX3400 config is not automatically captured, so switch
rebuild is slow and manual (gap, QL-REC09).
\item \textbf{RKE2:} rebuild rests on etcd quorum; etcd is not separately snapshotted
beyond VM-level backup.
\item \textbf{Git and docs:} GitHub plus local clones; the repos are the source of truth
for config-as-code.
\end{itemize}

\subsection{Recommendations}

Add a second, independent backup target (QL-REC03): the planned R730xd running its own
PBS with PBS-to-PBS remote sync, plus an offsite copy to Backblaze B2 for the datastore
and \code{bulk/fernanda}. Add the pve1 vzdump job (QL-REC04). See
\code{diagrams/recovery-dependencies.mmd}.
